\section{Development Lifecycle}\label{sec:lifecycle}

The development of \textsc{A3} followed a structured iterative methodology whose phases---\textbf{theorise $\to$ code $\to$ test $\to$ refine code $\to$ refine theory}---can be understood as a fixed-point computation over a product domain of formal specifications and executable implementations.
We describe the methodology here because it explains not only how the system was built but \emph{why} it has the layered architecture it does: each layer corresponds to a theoretical commitment that survived repeated refinement against executable semantics.

\subsection{Three Foundational Eras}

The project passed through three eras, each contributing a distinct formal layer to the final architecture.

\paragraph{Era~I: Quantitative Safety.}
The initial formulation modelled program safety as a metric problem.
Let $\mathcal{T}$ denote the set of program traces and $\mathcal{S} \subseteq \mathcal{T}$ the set of safe traces.
Rather than asking the Boolean question $\tau \in \mathcal{S}$, the system computed a distance $d(\tau, \partial\mathcal{S})$ from each trace to the safety boundary, treating this distance as a first-class semantic quantity amenable to optimisation.
Three principles survived to the final system:
(1)~safety admits continuous gradation---uncertainty has geometric structure;
(2)~quantitative semantics induces a natural priority ordering on analysis targets;
(3)~computability of $d(\tau, \partial\mathcal{S})$ enables guided repair.
What did not survive was the sufficiency assumption: in practice, the dominant engineering problem is not computing one deep metric but eliminating large volumes of false positives while preserving recall.

\paragraph{Era~II: Barrier-Certificate Separation.}
The system shifted from metric computation to set separation.
Let $\mathcal{I}$ denote the set of initial program states and $\mathcal{U}$ the set of unsafe states.
The core question became: does there exist a function $B : \Sigma \to \mathbb{R}$ over the state space $\Sigma$ satisfying
\[
  B(\mathbf{x}) \leq 0 \;\;\forall \mathbf{x} \in \mathcal{I}, \qquad
  B(\mathbf{x}) > 0 \;\;\forall \mathbf{x} \in \mathcal{U}, \qquad
  B(\mathbf{x}') \leq B(\mathbf{x}) \;\;\text{for every transition } \mathbf{x} \to \mathbf{x}'?
\]
If such a $B$ exists, the initial and unsafe regions are provably disjoint under the system's transition relation, and the corresponding bug report is a false positive.
This formulation, adapted from Prajna et al.~\cite{prajna2004safety,prajna2007framework}, provides the theoretical foundation for \textsc{A3}'s false-positive elimination engine.
The practical contribution of this era was recognising that barrier synthesis, when composed with SMT feasibility checks and CEGAR-style refinement, yields a high-throughput filter rather than an isolated proof technique.

\paragraph{Era~III: Executable Python Semantics.}
Instantiating the barrier framework over Python required defining the state space $\Sigma$, the transition relation $\to$, and the unsafe predicate $\mathcal{U}$ with respect to CPython's actual execution model.
This meant modelling bytecode-level control flow (including the exception table introduced in Python~3.11), normal and exceptional edges, frame and stack state, dynamic dispatch, and unknown library behaviour.
The transition relation is extracted directly from compiled bytecode via CFG construction (\S\ref{sec:engine}); the unsafe predicate is defined per bug type (e.g., for \texttt{DIV\_ZERO}: $\mathcal{U} = \{\sigma \in \Sigma \mid \sigma \text{ reaches a } \texttt{BINARY\_OP(/)} \text{ with divisor} = 0\}$).
Unknown callees are modelled via conservative contracts that overapproximate their possible return values and side effects.

\subsection{The Five-Phase Iteration}

The three eras did not proceed sequentially to completion; rather, they interleaved through a five-phase loop that constitutes the project's core methodology.
Let $T_i$ denote the formal theory and $A_i$ the analyser implementation at iteration $i$.
Each iteration produces $(T_{i+1}, A_{i+1})$ via:

\paragraph{Phase 1: Theorise.}
An LLM generates candidate formal hypotheses $\{h_1, \ldots, h_k\}$: new abstraction domains, candidate proof templates, verification-technique compositions, or cross-domain analogies (e.g., applying Lyapunov stability arguments to program-state reachability).
These hypotheses expand the search space of possible theories.
The generation is deliberately broad; soundness checking is deferred to subsequent phases.

\paragraph{Phase 2: Code.}
Each hypothesis $h_j$ is encoded as a modification to the analyser $A_i$, producing a candidate $A_i'$.
The encoding is constrained to the analyser's existing abstraction boundaries: crash summaries, barrier conditions, guard predicates, library contracts, and verification-engine routing rules.
This phase is deliberately mechanical---it translates formal claims into executable checks without introducing architectural novelty.

\paragraph{Phase 3: Test.}
The candidate $A_i'$ is evaluated against three increasingly expensive oracle sets:
(i)~a synthetic benchmark of 43~programs with known ground-truth labels;
(ii)~the BugsInPy corpus of 493~real bugs with buggy/fixed commit pairs;
(iii)~large-repository scans (e.g., DeepSpeed, 150\,k~LoC) to detect emergent effects---false-positive floods, performance regressions, confidence-score drift---invisible in unit tests.
A hypothesis that passes coding but fails testing is rejected: it is formally vacuous, operationally harmful, or both.

\paragraph{Phase 4: Refine Code.}
When $A_i'$ fails testing, the implementation is adjusted while holding the theory fixed.
Typical failure modes include: path explosion from unbounded symbolic exploration, over-conservative modelling of unknown callees that produces false positives, context loss across interprocedural boundaries, and duplicate findings from redundant proof attempts.
Refinements are kept local to the patch to bound the blast radius of each iteration.

\paragraph{Phase 5: Refine Theory.}
The most consequential phase.
When the code stabilises, the question is not ``does $A_{i+1}$ work?'' but ``what theory $T_{i+1}$ does $A_{i+1}$ actually implement?''
The formal specification is updated to match the surviving implementation: definitions are tightened, assumptions are made explicit, proof obligations are decomposed by semantic layer, and the conditions under which each barrier pattern $B_j$ is sound are precisely stated.
The revised theory $T_{i+1}$ then seeds the next iteration's hypothesis generation.

This loop converges in the sense that the gap $\|T_i - A_i\|$---informally, the discrepancy between what the theory promises and what the code delivers---decreases monotonically, because each iteration either corrects the code to match the theory (Phase~4) or corrects the theory to match the code (Phase~5).

\subsection{Illustration: Assertion-Escape Analysis}

Consider the property ``a failing assertion propagates as an uncaught exception.''
The initial theory $T_0$ modelled this as simple reachability: if a state $\sigma$ reaches an \texttt{RAISE\_VARARGS} instruction corresponding to \texttt{AssertionError}, the program is unsafe.
Testing revealed that this over-reports, because the theory neglected:
(i)~enclosing \texttt{except} handlers that catch \texttt{AssertionError};
(ii)~callers that wrap the call in \texttt{try}/\texttt{except};
(iii)~the \texttt{-O} flag, which disables assertions entirely;
(iv)~\texttt{finally} blocks that alter exception propagation.

Phase~4 refined the code to track exception-handler scopes via the bytecode exception table.
Phase~5 refined the theory: the unsafe predicate was narrowed from ``assertion raises'' to ``assertion raises \emph{and} no handler on any enclosing exception-table entry or caller-chain frame catches \texttt{AssertionError} or a superclass thereof.''
The same pattern---theory proposes, code implements, testing falsifies, refinement corrects both---repeated for unknown library calls, where the converged solution was contract-based overapproximation validated by concolic witness generation.

\subsection{Soundness Stratification}

The lifecycle produces a natural stratification of the analyser's claims by soundness level:

\begin{enumerate}
  \item \textbf{Theoretical soundness.}  AI-generated hypotheses are subjected to explicit proof obligations before encoding.  A hypothesis that cannot be stated as a predicate over $\Sigma$ with a decidable or semi-decidable checking procedure is rejected at Phase~1.
  \item \textbf{Implementation soundness.}  Analyser claims are verified against ground-truth oracles (synthetic benchmarks, BugsInPy commit pairs) and large-scale regression scans.  The BugsInPy self-improvement protocol (\S\ref{sec:bugsinpy}) and commit-level mining (\S\ref{sec:continuous}) provide continuous empirical validation.
  \item \textbf{Operational soundness.}  Residual findings that resist static proof and disproof are triaged by an agentic LLM (\S\ref{sec:triage}) whose classifications are bounded by the static pipeline's prior filtering: the LLM sees only the 1--3\% of candidates that survive deterministic elimination.  Baseline ratcheting (\S\ref{sec:ci}) enforces monotonic improvement of the deployed finding set.
\end{enumerate}

The asymmetry across these levels is deliberate: upstream AI \emph{expands} the hypothesis space; midstream formal machinery \emph{contracts} it via sound overapproximation; downstream agentic AI resolves only the residual where both proof and refutation are incomplete.
This separation prevents the generative phase from propagating noise into the verification phase, which is the mechanism by which the AI-slop feedback loop is inverted.
